\chapter{Bibliografía y recursos consultados}
A continuación se relacionan las principales fuentes de documentación que se han utilizado para la realización de esta práctica:

\begin{itemize}
    \item \textbf{OpenSSL Project.} \emph{OpenSSL Documentation: x509, req, genrsa y s\_client}. Disponible en \url{https://docs.openssl.org/}.
    \item \textbf{Apache HTTP Server Project.} \emph{Apache HTTP Server Version 2.4 Documentation — SSL/TLS Strong Encryption: How-To}. Disponible en \url{https://httpd.apache.org/docs/2.4/ssl/ssl_howto.html}.
    \item \textbf{Apache HTTP Server Project.} \emph{Module mod\_ssl}. Disponible en \url{https://httpd.apache.org/docs/2.4/mod/mod_ssl.html}.
    \item \textbf{Debian/Kali Linux.} \emph{Manual de \texttt{update-ca-certificates(8)}}. Disponible mediante \texttt{man update-ca-certificates} y en \url{https://manpages.debian.org/}.
    \item \textbf{Mozilla Developer Network (MDN).} \emph{Transport Layer Security (TLS)}. Disponible en \url{https://developer.mozilla.org/en-US/docs/Web/Security/Transport_Layer_Security}.
    \item \textbf{Asistente generativo:} se ha utilizado \emph{Claude} (Anthropic), disponible en \url{https://claude.ai}, para resolver dudas puntuales de sintaxis de \texttt{openssl} y de configuración de Apache, conforme a lo que se permite en el enunciado de la práctica.
\end{itemize}
